%% kapa: exam=1 round=25 year=2014 period=1 form=A subject=economics subject_ko=경제원론 seq=25 number=65
%% source: Q-Net 감정평가사 기출문제 (cst00309) artlSeq=5101207 file=extracted/1008/5101207/2014년 제25회 감정평가사 1차시험 기출문제/1교시(A형).hwp
%% conversion: kordoc 4.24.51 -> tools/gen1.py (auto); review=proofread-sonnet; flags=none
\begin{question}{65}
다음 ( )안의 내용으로 알맞은 것은?

\begin{framedtext}
원유수입가격 상승시 원유수입국의 소비자물가지수는 ( ㄱ )하고, 생산자물가지수는 ( ㄴ )하며, GDP 디플레이터는 ( ㄷ )한다. \par
\end{framedtext}
\begin{choices}
  \item ㄱ: 불변, ㄴ: 불변, ㄷ: 상승
  \item ㄱ: 상승, ㄴ: 상승, ㄷ: 불변
  \item ㄱ: 상승, ㄴ: 불변, ㄷ: 상승
  \item ㄱ: 불변, ㄴ: 상승, ㄷ: 상승
  \item ㄱ: 상승, ㄴ: 상승, ㄷ: 상승
\end{choices}
\end{question}
\answer{5}
